\chapter*{摘要}
\addcontentsline{toc}{chapter}{摘要}

随着文本到语音合成（TTS，Text-to-Speech）、语音克隆和大语言模型（LLM，Large
Language
Model）驱动的生成式语音技术快速发展，真实语音样本被未授权采集、编码和复刻的技术门槛不断降低。公开视频、直播录屏、播客节目、有声读物、会议录音、企业客服语音和宣传音视频中均包含大量可识别的个人或组织声音资产。攻击者一旦获取这些公开语音，就可能利用零样本语音克隆、少样本微调式语音克隆或基于语音分词器（Speech
Tokenizer）的新型语音合成系统，生成与目标说话人高度相似的伪造语音，进而引发身份冒用、电信诈骗、虚假内容传播、企业声音仿冒和个人声音权益滥用等安全风险。声音作为一种具有身份属性、传播频繁且难以更换的数字资产，已经成为生成式人工智能安全治理中的重要保护对象。

现有语音安全治理手段多集中于事后检测、水印溯源和内容审核环节。这类方法能够在合成语音已经产生或传播之后辅助判断其真实性，也可以为平台审核和事后追责提供依据，但难以阻止攻击者在语音公开发布后收集原始声音样本。一旦真实语音被大规模采集，攻击者便可能在检测系统介入之前完成参考音频整理、语义标注、说话人特征提取和语音克隆建模。因此，仅依赖事后检测难以充分解决声音资产源头泄露问题。

针对上述不足，本作品设计并实现了一套面向声音资产安全的源头主动防护系统
V-Guard，面向公开语音发布、音视频内容上传、企业语音素材外发和反诈安全演示等场景，在语音发布前为原始音频生成细粒度保护扰动，输出受保护音频，从源头降低未授权语音克隆系统对原始声音资产的稳定提取能力。

V-Guard的设计基于对语音克隆系统信息依赖结构的分析。语音克隆至少需要从参考音频中提取两类信息：一类是语义表示链路，即机器提取语音内容、发音结构和语义token的过程，回答``说了什么''；另一类是声音身份表示链路，即机器提取说话人身份、音色和声学条件的过程，回答``是谁在说、怎么说''。系统通过语义代理模型组和声音身份代理模型组对上述两类链路分别建模，使保护扰动同时作用于两个表示空间。在此基础上引入心理声学约束控制扰动的可感知性，避免保护后音频出现明显听感失真。

\textbf{本作品的创新主要体现在四个方面：}（1）将防护位置从事后检测前移至语音公开发布之前，在源头降低语音样本对未授权克隆系统的利用价值；（2）构建语义与声音身份双链路联合防护框架，其中语义链路从传统的ASR文本输出层面下沉至语音编码器的中间表示层面，以更根本的方式破坏语音内容的可理解性；声音身份链路通过多编码器集成削弱已知和未知克隆模型对目标说话人身份的提取能力；（3）引入心理声学掩蔽模型约束扰动的时频分布，在防护强度与音频质量之间提供可调节的权衡空间；（4）将上述算法能力封装为包含接入、保护、评估和导出的完整网页端工程平台，使主动防护技术具备可操作性和可验证性。

对作品的评估实验方面，本作品基于 LibriTTS
测试样本构建``原始音频---保护音频''成对评估流程，并从语义保护、声音身份保护、音频质量、鲁棒性、消融实验和系统性能等角度进行验证。测试结果表明，受保护音频能够显著干扰机器语义识别链路，在
Whisper Medium 和 wav2vec2-base-960h 上均产生明显 WER
提升；在语音克隆验证中，以保护音频作为 XTTS-v2
参考音频生成的克隆语音，其 ECAPA-TDNN
说话人相似度较原始参考条件显著下降，说明系统能够削弱克隆后端对原始说话人身份的复刻能力。进一步的高保真参数组、跨克隆后端迁移性、信道变换鲁棒性和损失项消融实验也表明，V-Guard
具备一定参数可调性、迁移性和传播链路残留防护能力。

V-Guard
系统不能替代法律监管、平台审核、水印溯源和深度伪造检测，但可以作为上述治理手段之前的一道前置保护层，在声音样本被采集、编码和克隆利用之前降低其机器可利用性，从而为生成式语音时代的声音隐私保护、声音资产管理和反诈安全教育提供工程化支撑。

\textbf{关键词}：声音资产；主动防护；语音克隆；语义表示链路；声音身份表示链路；心理声学约束；受保护音频；语音安全
